\chapter{Expectation}

In the last chapter we cataloged named distributions like the binomial,
geometric, and Poisson. Each gives a full list of probabilities for every
possible outcome. But if you describe an uncertain setup to an engineer or an
investor, they rarely ask for the whole list. They ask one plain question: on
average, what should I expect?

That average is the \emph{expectation} (or \emph{expected value}, or
\emph{mean}). You can picture it in two complementary ways: as the long-run
average if you repeat an experiment many times, or as the physical balance
point where a plot of the probabilities sits level on a knife edge. After
working with both pictures, we will use linearity of expectation to find
averages for sums of dependent variables, and see where averages break down in
the St.\ Petersburg paradox.

\chapterline

\section{Two pictures of the average}
\label{ch08:two-pictures}

\subsection{The long-run-average picture}

Imagine a game. You roll a fair six-sided die, and you win in dollars whatever
number comes up. Some rounds you win \$1, some rounds \$6, most rounds something
in between. If you play many, many rounds, what is your average winning per
round?

Intuitively, each face shows up about one-sixth of the time, so in $6000$ rounds
you would see roughly $1000$ of each face. Your total winnings would be about
\[
  1000\cdot 1 + 1000\cdot 2 + \cdots + 1000 \cdot 6 = 1000\cdot 21 = 21000,
\]
and dividing by the $6000$ rounds gives an average of $21000/6000 = 3.5$ dollars
per round. The $1000$'s cancelled, and we were left with
\[
  \frac{1+2+3+4+5+6}{6} = \frac{21}{6} = 3.5.
\]
The fraction of rounds with each face -- the probability -- is what survives.
That is the long-run-average picture. The expectation is the per-round payoff you
approach as the number of rounds grows.

\begin{intuitionbox}
Expectation answers the question: \emph{if I could repeat this exactly the same
way an enormous number of times and average the results, what number would I land
on?} You do not need the experiment to actually be repeatable in real life. The
expectation is a property of the distribution, computed as if it were repeatable.
\end{intuitionbox}

\subsection{The center-of-mass picture}

Here is the second picture, the one that makes expectation feel
geometric and inevitable. Draw the probability mass function of your random
variable as little weights sitting on a number line: at each possible value $x$
you place a lump of mass equal to $\P(X = x)$. The total mass is $1$ (the
probabilities sum to one). Now ask a physics question: where is the balance point
of this arrangement of weights?

\begin{figure}[H]
\centering
\begin{tikzpicture}
\begin{axis}[
  width=12cm, height=5.2cm,
  title style={font=\small},
  title={Expectation as the Physical Balance Point of the Die PMF},
  ybar, bar width=14pt,
  xlabel={value $x$}, ylabel={$\P(X=x)$},
  xtick={1,2,3,4,5,6}, ymin=0, ymax=0.25, clip=false,
  ytick={0.1,0.2},
  enlarge x limits=0.12,
  axis lines=left,
]
\addplot[fill=distblue!55,draw=distblue] coordinates {(1,0.1667)(2,0.1667)(3,0.1667)(4,0.1667)(5,0.1667)(6,0.1667)};
\draw[densityred,very thick,->] (axis cs:3.5,0) -- (axis cs:3.5,0.205);
\node[densityred,anchor=south] at (axis cs:3.5,0.205) {\small balance point $=\E X = 3.5$};
\end{axis}
\draw[thick] (0.6,-1.05) -- (10.4,-1.05);
\fill[black] (5.21,-1.05) -- (4.99,-1.32) -- (5.43,-1.32) -- cycle;
\node at (5.21,-1.58) {\small fulcrum at $3.5$: the weights balance};
\end{tikzpicture}
\caption{The fair die as a center of mass. Six equal weights sit at
$1,\dots,6$; the arrangement balances exactly at $3.5$, which is the
expectation. The fulcrum sits under the mean, not under any actual outcome.}
\label{ch08:fig-die-balance}
\end{figure}

For the die, six equal weights at $1$ through $6$ balance symmetrically at the middle,
$3.5$. The expectation \emph{is} the balance point. And notice that the
balance point need not be a value the variable can ever take. You never roll a
$3.5$. The expected value is not the ``most likely'' value, and it need not even
be a \emph{possible} value. It is the center of mass of the probabilities.

\begin{recoverybox}
\textbf{Pause and think.} Suppose a game pays \$0 with probability $0.99$ and
\$1000 with probability $0.01$. Where is the center of mass? It balances at
$0\cdot 0.99 + 1000\cdot 0.01 = 10$ dollars. Notice that $10$ is nowhere near
either possible outcome ($0$ or $1000$). The mean does not describe what
happens most frequently; it marks the physical balance point of the entire
distribution.
\end{recoverybox}

\begin{historybox}
The concept of expectation is older than the modern mathematical definition of
probability itself. In the summer of 1654, a French nobleman and seasoned
gambler, Antoine Gombaud (the Chevalier de M\'er\'e), posed a classic puzzle to
Blaise Pascal: the \emph{problem of points}. Two players put equal stakes into a
pot for a game of chance, agreeing that whoever first wins a set number of
rounds takes the entire pot. If the game must be abandoned early when one
player leads, how should the pot be divided fairly?

For over a century, Italian scholars including Luca Pacioli (in 1494) and
Niccol\`o Tartaglia (in 1556) had argued over rival rules, dividing the pot in
proportion to rounds already won or rounds remaining. Their methods led to
absurd contradictions -- if a game to ten wins were stopped after one round at
$1\text{--}0$, Pacioli's rule awarded the entire pot to the leader! The
mathematicians were stuck because they lacked any formal framework for
probability; they had only proportional arithmetic.

Pascal initiated a famous exchange of letters with Pierre de Fermat. Their
breakthrough was to ignore the past and look strictly at the future: evaluate
each player's fair share as the average payoff they could \emph{expect} to
receive across all possible ways the remaining rounds could unfold. In 1657, the
Dutch scientist Christiaan Huygens learned of their correspondence and published
the first formal probability treatise, \emph{De Ratiociniis in Ludo Aleae}
(``On Reasoning in Games of Chance''). Huygens took as his very first
proposition the calculation of \emph{expectatio} -- the fair price of an
uncertain gamble. Historically, mathematicians understood what a fair price was
before they had a settled definition of probability.

Yet this foundational concept soon met controversy. In 1713, Nicolaus Bernoulli
proposed a coin-tossing game (now called the St.\ Petersburg paradox) whose
calculated expectation was infinite, even though sensible gamblers would not pay
even twenty dollars to play. This forced his cousin Daniel Bernoulli in 1738 to
propose \emph{expected utility} -- the insight that a dollar is worth less to a
wealthy person than to a pauper, so people maximize subjective satisfaction
rather than raw expected money. Even today, understanding how people make
decisions when faced with tiny probabilities of catastrophic losses or gigantic
jackpots remains an open, contested frontier at the intersection of probability,
economics, and cognitive psychology.
\end{historybox}

\subsection{The definition}

We can now write down the definition that both pictures point to. For a discrete
random variable $X$ taking values $x_1, x_2, \dots$ with probabilities
$p_i = \P(X = x_i)$, weight each value by its probability and add:

\begin{definition}[Expectation, discrete case]
\label{ch08:def-expectation}
Let $X$ be a discrete random variable with possible values $x_1,x_2,\dots$. Its
\emph{expectation} (or \emph{expected value}, or \emph{mean}) is
\[
  \E X \;=\; \sum_i x_i\,\P(X = x_i),
\]
provided the sum converges absolutely, i.e.\ $\sum_i \abs{x_i}\,\P(X=x_i) < \infty$.
We also write $\mu$ or $\mu_X$ for $\E X$.
\end{definition}

The absolute-convergence proviso is not pedantry: if the positive and negative
parts both blow up, no single number deserves to be called ``the average,'' and
rearranging the sum could give you anything. We will meet a clean example (St.\
Petersburg) where the sum genuinely diverges and the expectation is $+\infty$.

\begin{focusbox}
Both pictures are the \emph{same formula}, read two ways.
\begin{itemize}
  \item \textbf{Long-run average:} $\P(X=x_i)$ is the fraction of trials landing
  at $x_i$, and $\sum_i x_i\,\P(X=x_i)$ is the limiting average payoff.
  \item \textbf{Center of mass:} $\P(X=x_i)$ is the mass at position $x_i$, and
  $\sum_i x_i\,\P(X=x_i)$ is the formula for the balance point (in physics and
  statistics, the \emph{first moment}, meaning distance weighted by mass, which
  we explore fully in Chapter~9).
\end{itemize}
\end{focusbox}

\begin{example}[A loaded bet]
\label{ch08:ex-loaded-bet}
A carnival game: you pay nothing to play. With probability $0.7$ you win \$1,
with probability $0.25$ you win \$5, and with probability $0.05$ you lose \$20.
Let $X$ be your net gain. Then
\[
  \E X = (1)(0.7) + (5)(0.25) + (-20)(0.05)
       = 0.7 + 1.25 - 1.0 = 0.95.
\]
On average you make $95$ cents per play. Most plays you win a little; the rare
\$20 loss is real but does not eat the whole edge. If the carnival ran this game
it would slowly go broke -- which is exactly why real carnival games have
$\E X < 0$ for the player.
\end{example}

\begin{example}[A constant is its own mean]
\label{ch08:ex-constant}
If $X = c$ always (a degenerate ``random'' variable), then $\E X = c\cdot 1 = c$.
Obvious, but worth stating: averaging something that never varies gives back the
thing. This is the base case that linearity will lean on.
\end{example}

\subsection{First encounters with Chapter 7 distributions}

Armed with the definition, we can immediately compute the mean for some of the
named distributions we introduced in Chapter~7.

\begin{example}[Mean of a Bernoulli trial]
\label{ch08:ex-bern-mean}
Let $X \sim \Bern(p)$, so $X$ takes value $1$ (success) with probability $p$ and
value $0$ (failure) with probability $1-p$. Directly from the definition:
\[
  \E X = 0\cdot (1-p) + 1\cdot p = p.
\]
The expected value of a single trial is simply its probability of success.
\end{example}

\begin{example}[Mean of a Poisson variable]
\label{ch08:ex-poisson-mean}
Let $X \sim \Pois(\lambda)$ count rare events over an interval, with PMF
$\P(X = k) = e^{-\lambda}\lambda^k / k!$ for $k = 0, 1, 2, \dots$.
Computing the mean directly from the definition:
\[
  \E X = \sum_{k=0}^\infty k\,\frac{e^{-\lambda}\lambda^k}{k!}
       = \sum_{k=1}^\infty \frac{e^{-\lambda}\lambda^k}{(k-1)!}
       = \lambda e^{-\lambda} \sum_{k=1}^\infty \frac{\lambda^{k-1}}{(k-1)!}
       = \lambda e^{-\lambda} e^{\lambda} = \lambda.
\]
The parameter $\lambda$ introduced in Chapter~7 as an arrival rate
equals the expected number of events.
\end{example}

What about the other distributions from Chapter~7? For the binomial,
summing $\sum_{k=0}^n k \binom{n}{k} p^k (1-p)^{n-k}$ requires binomial identities,
and for the hypergeometric distribution (cards dealt without replacement), summing
the PMF directly is even harder. We need a method that bypasses algebraic sums.
That method begins with a direct connection between probabilities and averages.

\section{Expectation of an indicator is a probability}
\label{ch08:indicators}

Before linearity, we need one basic identity that underlies many arguments
in probability. Recall the \emph{indicator} random variable of an event $A$:
\[
  \ind_A = \begin{cases} 1 & \text{if } A \text{ occurs},\\ 0 & \text{otherwise.}\end{cases}
\]
It is the simplest non-constant random variable there is: it just flags whether
$A$ happened. What is its expectation? By the definition, $\ind_A$ takes value $1$
with probability $\P(A)$ and value $0$ with probability $1-\P(A)$, so
\[
  \E[\ind_A] = 1\cdot \P(A) + 0\cdot \bigl(1-\P(A)\bigr) = \P(A).
\]

\begin{formalbox}
\textbf{The bridge.} For any event $A$,
\[
  \boxed{\;\E[\ind_A] = \P(A).\;}
\]
Probability is a special case of expectation. Every probability you have ever
computed is already an expected value -- the expected value of a yes/no flag.
\end{formalbox}

This little identity is the reason expectation is more fundamental than it first
appears. It lets us \emph{translate counting questions into averaging questions},
and averaging questions are precisely where the next tool proves exceptionally powerful.

\begin{recoverybox}
\textbf{Pause and think.} Roll a fair six-sided die, and let $A$ be the event
that the roll is even ($A = \{2, 4, 6\}$). What is $\E[\ind_A]$? By our bridge
identity, $\E[\ind_A] = \P(A) = 3/6 = 1/2$. If you tally a $1$ whenever an even
roll occurs and $0$ otherwise, the long-run average of your tally is precisely
the probability of rolling an even number.
\end{recoverybox}

\begin{happeningbox}
Why is this so useful? Because we are about to learn that expectation splits over
sums, even when the pieces depend on each other. So if you can write
the quantity you care about as a \emph{sum of indicators} -- ``how many of these
events happened?'' -- then its expectation is just the \emph{sum of the
probabilities}, $\E\bigl[\sum_i \ind_{A_i}\bigr] = \sum_i \P(A_i)$. A hard
combinatorial count collapses into an easy addition of probabilities. Keep this
phrase in mind: \emph{count by summing indicators}.
\end{happeningbox}

\section{Linearity of expectation}
\label{ch08:linearity}

Here is the theorem that earns its place as the single most useful fact in this
book.

\begin{theorem}[Linearity of expectation]
\label{ch08:thm-linearity}
Let $X$ and $Y$ be random variables on the same probability space with finite
expectations, and let $a, b$ be constants. Then
\[
  \E[aX + bY] = a\,\E X + b\,\E Y.
\]
More generally, for any random variables $X_1,\dots,X_n$ with finite means,
\[
  \E\!\left[\sum_{i=1}^n X_i\right] = \sum_{i=1}^n \E X_i.
\]
\emph{No independence is required.} The $X_i$ may depend on each other in any way
whatsoever.
\end{theorem}

Read that last sentence twice. The expectation of a sum equals the sum of
expectations in all cases, regardless of whether the variables are independent.

\begin{whybox}
Why is linearity true, with no independence assumption? Go back to the most
basic level: a random variable is a function $X(\omega)$ of the underlying
outcome $\omega$, and its expectation is an average over outcomes weighted by
their probabilities:
\[
  \E X = \sum_{\omega} X(\omega)\,\P(\omega).
\]
Now for the sum $X + Y$, at \emph{each individual outcome} $\omega$ we simply add
the two numbers: $(X+Y)(\omega) = X(\omega) + Y(\omega)$. There is no interaction
-- addition of two numbers is done pointwise. So
\[
  \E[X+Y] = \sum_\omega \bigl(X(\omega) + Y(\omega)\bigr)\P(\omega)
          = \sum_\omega X(\omega)\P(\omega) + \sum_\omega Y(\omega)\P(\omega)
          = \E X + \E Y.
\]
Independence never entered because we never had to split a probability of a
\emph{joint} event. We only rearranged a sum -- and sums of numbers can always be
regrouped. Dependence is about how $X$ and $Y$ \emph{co-occur}, i.e.\ about joint
probabilities; but linearity never touches joint probabilities. It just adds two
real numbers at each outcome and then totals up. \emph{That} is why it does not
care about dependence.
\end{whybox}

\begin{trapbox}
Linearity is about \emph{sums}, not products. In general
\[
  \E[XY] \ne (\E X)(\E Y),
\]
and the discrepancy is exactly the covariance (next chapter). Equality of the
product holds when $X \indep Y$, but not otherwise. Likewise $\E[g(X)] \ne
g(\E X)$ for a nonlinear $g$ -- that is Jensen's inequality, also coming later.
Linearity is a statement about $aX+bY$, and only that.
\end{trapbox}

\subsection{Why linearity is the single most useful tool}

When you want to find $\E X$ directly from the definition, you need the full
distribution of $X$ -- every value and its probability. For complicated $X$ that
distribution can be a nightmare. Linearity offers an escape hatch: \emph{break
$X$ into a sum of simple pieces}, find the expectation of each piece, and
add. The component variables may be strongly dependent, yet no joint distribution
needs to be computed. You never have to find the distribution of $X$ as a whole.

The recipe is almost always the same three steps.

\begin{focusbox}
\textbf{The indicator method.}
\begin{enumerate}
  \item Write the quantity of interest as a sum of indicators:
        $X = \sum_i \ind_{A_i}$ (``count how many $A_i$ happened'').
  \item Use $\E[\ind_{A_i}] = \P(A_i)$ on each piece.
  \item Add up by linearity: $\E X = \sum_i \P(A_i)$.
\end{enumerate}
The dependence between the $A_i$ is irrelevant. Just sum the \emph{marginal
probabilities} -- the probability of each event $A_i$ on its own, regardless
of what happens to the others.
\end{focusbox}

\subsection{Warm-up: the mean of a binomial without any algebra}

How many heads do you expect in $n$ flips of a coin that lands heads with
probability $p$? Let $X \sim \Bin(n,p)$ be the number of heads. The direct
approach is
\[
  \E X = \sum_{k=0}^n k\binom{n}{k}p^k(1-p)^{n-k},
\]
a sum you \emph{can} grind out with a binomial-identity trick, but it is fiddly.
Now do it the easy way. Let $X_i = \ind\{\text{flip } i \text{ is heads}\}$. Then
\[
  X = X_1 + X_2 + \cdots + X_n,
\]
and each $X_i$ is an indicator with $\E X_i = \P(\text{heads on flip } i) = p$.
By linearity,
\[
  \E X = \sum_{i=1}^n \E X_i = \sum_{i=1}^n p = np.
\]
Three lines, no binomial coefficients. The flips here happen to be independent,
but -- as the next example illustrates -- that independence is never required.

\subsection{Dependent trials: the mean of a hypergeometric distribution}
\label{ch08:hypergeometric-mean}

Now see what happens when the trials are heavily \emph{dependent}. Recall from
Chapter~7 the hypergeometric distribution: you draw $n$ cards without
replacement from a deck of $N$ cards containing $K$ aces. Let $X$ be the number
of aces drawn.

If you attempt this with the PMF, you must evaluate the sum
\[
  \E X = \sum_{k} k \cdot \frac{\binom{K}{k}\binom{N-K}{n-k}}{\binom{N}{n}},
\]
which demands Vandermonde convolution identities. With the indicator method,
it takes the same three lines as the binomial.

Define $X_i = \ind\{\text{the } i\text{th card drawn is an ace}\}$. Then $X =
X_1 + \dots + X_n$. Before any cards are revealed, every draw is equally likely
to be any of the $N$ cards in the deck, so by symmetry
\[
  \P(X_i = 1) = \frac{K}{N}.
\]
These indicators are dependent: drawing an ace on the first card strictly
reduces the chance of an ace on the second card. Yet linearity does not care!
Summing the marginal expectations gives
\[
  \E X = \sum_{i=1}^n \E X_i = \sum_{i=1}^n \frac{K}{N} = n \cdot \frac{K}{N}.
\]
Setting $p = K/N$, the mean is simply $np$, identical to the binomial formula.
In Chapter~7 we emphasized the conceptual boundary between sampling with
replacement (independent, binomial) and sampling without replacement (dependent,
hypergeometric). For expected values, that boundary dissolves: linearity handles
both with the exact same ease.

\section{The astonishing power: hat-matching and dependence}
\label{ch08:hat-matching}

Now the showcase -- the example that tends to convert skeptics.

\begin{example}[The hat-check / matching problem]
\label{ch08:ex-hats}
$n$ people check their hats at a restaurant. A careless attendant returns the
hats in a completely random order -- a uniformly random permutation. On average,
how many people get their \emph{own} hat back?

Let $X$ be the number of people who receive their own hat. Finding the
distribution of $X$ is genuinely hard: it involves \emph{derangements} (the
number of permutations with no fixed point) and an inclusion--exclusion sum. But
we do not want the distribution. We want $\E X$, and linearity makes it immediate.

Define the indicator $X_i = \ind\{\text{person } i \text{ gets their own hat}\}$.
Then $X = X_1 + \cdots + X_n$. For each person, the hat they receive is equally
likely to be any of the $n$ hats, so
\[
  \E X_i = \P(\text{person } i \text{ gets own hat}) = \frac{1}{n}.
\]
By linearity,
\[
  \E X = \sum_{i=1}^n \E X_i = n \cdot \frac{1}{n} = 1.
\]
\emph{On average exactly one person gets their own hat back, no matter how large
$n$ is.} Ten people or ten million, the expected number of matches is $1$.
\end{example}

\begin{happeningbox}
Notice what we did not have to compute. The indicators $X_i$ are \emph{not} independent: if
the first $n-1$ people all got their own hats, then the last person is
\emph{forced} to get theirs too (only one hat remains). The events are
dependent. Yet linearity applies directly, because -- as the proof showed --
it only ever adds the marginal probabilities $\P(X_i = 1) = 1/n$, and the
marginal probability of person $i$'s match does not depend on anyone else's outcome. The
dependence would matter if we computed the \emph{variance} of $X$ (which
involves $\E[X_iX_j]$, a joint quantity) -- but it does not enter the mean.
\end{happeningbox}

\begin{figure}[H]
\centering
\begin{tikzpicture}[
  lane/.style={draw=#1!70, fill=#1!8, rounded corners=3pt, thick, inner sep=6pt},
  box/.style={draw=#1!80, fill=#1!18, rounded corners=2pt, font=\footnotesize, align=center, thick},
  arr/.style={->, >=stealth, thick, draw=black!70}
]
  % Lane 1: The Full Distribution Route (Hard)
  \node[lane=trapred, minimum width=12cm, minimum height=1.6cm] (lane1) at (0, 1.4) {};
  \node[anchor=north west, font=\footnotesize\bfseries\color{trapred!80!black}] at (lane1.north west)
    {Route A: The Full PMF Route (Combinatorially Intractable)};
  \node[box=trapred] (stepA1) at (-4.2, 1.25) {Permutation space\\$n!$ outcomes};
  \node[box=trapred] (stepA2) at (-0.2, 1.25) {Derangements $!n$ \&\\Inclusion--Exclusion};
  \node[box=trapred] (stepA3) at (4.0, 1.25) {Full PMF $\P(X=k)$\\$\sum k \cdot \P(X=k)$};
  \draw[arr, trapred!80] (stepA1) -- (stepA2);
  \draw[arr, trapred!80] (stepA2) -- (stepA3);

  % Lane 2: The Indicator Method (Fast & Clean)
  \node[lane=targetgreen, minimum width=12cm, minimum height=1.6cm] (lane2) at (0, -0.7) {};
  \node[anchor=north west, font=\footnotesize\bfseries\color{targetgreen!80!black}] at (lane2.north west)
    {Route B: Linearity of Indicators (Clean Decomposition)};
  \node[box=targetgreen] (stepB1) at (-4.2, -0.85) {Atomic Indicators\\$X = \sum_{i=1}^n \ind_{A_i}$};
  \node[box=targetgreen] (stepB2) at (-0.2, -0.85) {Marginal Chances\\$\E[\ind_{A_i}] = \P(A_i) = \frac{1}{n}$};
  \node[box=targetgreen] (stepB3) at (4.0, -0.85) {Linearity of $\E$\\$\sum_{i=1}^n \frac{1}{n} = 1$};
  \draw[arr, targetgreen!80] (stepB1) -- (stepB2);
  \draw[arr, targetgreen!80] (stepB2) -- (stepB3);

  % Final Result
  \node[draw=sourceblue!80, fill=sourceblue!15, rounded corners=2pt, font=\small\bfseries, align=center]
    (res) at (0, -2.3) {Expected Matches: $\E[X] = 1$ (for any population size $n \ge 1$)};
  \draw[arr, trapred!80] (stepA3.east) to[out=0, in=10] (res.east);
  \draw[arr, targetgreen!80] (stepB3.south) -- (stepB3.south |- res.north);
\end{tikzpicture}
\caption{Why linearity bypasses large combinatorial sums. To find $\E[X]$, Route~A constructs the full distribution through derangements and inclusion--exclusion. Route~B slices $X$ into indicators $\ind_{A_i}$ and sums their marginal expectations. The dependencies between people never need to be untangled.}
\label{ch08:fig-indicator-bypass}
\end{figure}

\begin{labbox}
Before reading on, try this with the same one-line method. You shuffle a standard
deck and deal the cards face up one at a time while reciting ``ace, two, three,
\dots, king, ace, two, \dots'' through the $52$ cards in a fixed order. A
``match'' is when the card you name equals the card you flip. What is the expected
number of matches? \emph{Hint:} indicator per position, each matches with
probability... what? \\[2pt]
\emph{Answer:} each of the $52$ positions matches its named rank with probability
$\tfrac{4}{52}=\tfrac1{13}$ (four cards of the named rank among $52$), so the
expected number of matches is $52\cdot\tfrac1{13} = 4$. Dependence between
positions is irrelevant.
\end{labbox}

\begin{example}[Expected number of records]
\label{ch08:ex-records}
A sequence of $n$ distinct numbers arrives in random order (all $n!$ orders
equally likely). Position $i$ is a \emph{record} if its value exceeds everything
before it. How many records do you expect?

Let $R_i = \ind\{\text{position } i \text{ is a record}\}$. Position $i$ is a
record exactly when the largest of the first $i$ values happens to sit in slot
$i$. By symmetry every one of those $i$ values is equally likely to be the
largest, so $\P(R_i = 1) = 1/i$. Hence
\[
  \E[\text{records}] = \sum_{i=1}^n \frac1i = H_n \approx \ln n + 0.577.
\]
The expected number of records grows only \emph{logarithmically}. After a
thousand observations you expect about $\ln 1000 + 0.58 \approx 7.5$ record-breaking
values. Here $H_n = \sum_{k=1}^n \tfrac1k = 1 + \tfrac12 + \tfrac13 + \dots + \tfrac1n$
is the $n$th \emph{harmonic number} -- the sum of the first $n$ unit fractions.
It grows logarithmically, tracking $\ln n$ plus the Euler--Mascheroni constant
$\gamma \approx 0.5772$. This same $H_n$ governs the coupon-collector problem,
which we solve in Section~\ref{ch08:gallery}.
\end{example}

\begin{figure}[H]
\centering
\begin{tikzpicture}
\begin{axis}[
  width=12cm, height=5.6cm,
  title style={font=\small},
  title={Indicator Method Scaling: Matches vs.\ Harmonic Growth of Records},
  xlabel={number of people / observations $n$},
  ylabel={expected count},
  domain=1:50, samples=50,
  legend pos=north west, legend cell align=left,
  axis lines=left, ymin=0, ymax=5,
  xmin=1, xmax=50,
]
\addplot[densityred,very thick,domain=1:50] {1};
\addlegendentry{hat matches: constant $1$}
\addplot[distblue,very thick,domain=1:50] {ln(x)+0.5772};
\addlegendentry{records: $H_n \approx \ln n + 0.577$}
\end{axis}
\end{tikzpicture}
\caption{Two indicator-method results. The expected number of self-matched hats
stays pinned at $1$ for every $n$, while the expected number of records grows
logarithmically as $\ln n$. Both are direct linearity computations over dependent
indicators.}
\label{ch08:fig-hats-records}
\end{figure}

\begin{applicationbox}
\textbf{Linearity in computer science and data science.} The indicator method is
widely used in \emph{randomized algorithm analysis}.
\begin{itemize}
  \item \textbf{Hashing / load balancing.} Throw $m$ items into $n$ hash buckets
  uniformly. The expected number of items in a given bucket is, by indicators,
  $m\cdot \tfrac1n = m/n$ -- and the expected number of \emph{empty} buckets is
  $n\,(1-\tfrac1n)^m \approx n\,e^{-m/n}$.
  \item \textbf{Expected comparisons in QuickSort.} Let $A_{ij}$ be the event
  that the $i$th and $j$th smallest elements are ever compared. One shows
  $\P(A_{ij}) = \tfrac{2}{j-i+1}$, and summing these (highly dependent!)
  indicators gives the famous $\E[\#\text{comparisons}] \approx 2n\ln n$.
  \item \textbf{Count-distinct sketches} (HyperLogLog and friends) are built on
  expectations of indicator-like statistics over hashed streams.
\end{itemize}
In each case nobody computes a full distribution. They sum marginal
probabilities, because expectation lets them.
\end{applicationbox}

\section{LOTUS: the expectation of a function}
\label{ch08:lotus}

So far we have averaged $X$ itself. But usually we care about some
\emph{function} of $X$ -- your payoff is $g(X)$, your loss is $(X-\mu)^2$, your
utility is $\sqrt{X}$. How do you compute $\E[g(X)]$?

The tempting wrong move is to first find the distribution of the new variable
$Y = g(X)$, then apply the definition to $Y$. That works but is usually painful:
deriving the distribution of $g(X)$ can be hard, and it is wasted effort. There
is a shortcut so handy it has a nickname.

\begin{theorem}[Law of the Unconscious Statistician, discrete case]
\label{ch08:thm-lotus}
Let $X$ be a discrete random variable with pmf $p_X$, and let $g$ be a function.
Then
\[
  \E[g(X)] \;=\; \sum_{x} g(x)\,\P(X = x),
\]
provided $\sum_x \abs{g(x)}\,\P(X=x) < \infty$.
\end{theorem}

In words: to average $g(X)$, just apply $g$ to each value of $X$ and weight by the
\emph{original} probabilities of $X$. You never have to find the distribution of
$g(X)$ at all. The nickname pokes fun at the statistician who uses this formula
``unconsciously,'' as if it were the definition -- but it is a theorem, and it
needs a (short) proof.

\begin{whybox}
Why is LOTUS true? The real worry is that $g$ might collapse several different
$x$-values onto the same $y$-value, so the probabilities should pile up. They do,
and the bookkeeping works out. Group the $x$'s by the value $y=g(x)$ they produce:
\begin{align*}
  \E[g(X)] &= \sum_y y\,\P\bigl(g(X)=y\bigr) && \text{(definition, applied to } Y=g(X))\\
           &= \sum_y y \sum_{x:\,g(x)=y} \P(X=x) && \text{(the event } g(X)=y \text{ is the union of those } x)\\
           &= \sum_y \sum_{x:\,g(x)=y} g(x)\,\P(X=x) && \text{(inside this group } y=g(x))\\
           &= \sum_x g(x)\,\P(X=x). && \text{(the double sum reassembles all } x)
\end{align*}
The collapsing is handled automatically: every $x$ contributes $g(x)\P(X=x)$
exactly once, whether or not it shares its $g$-value with other $x$'s.
\end{whybox}

\begin{connectionbox}
LOTUS is the discrete cousin of a change-of-variables you already know from
calculus. In the continuous chapters $\E[g(X)] = \int_{-\infty}^{\infty}
g(x)\,f_X(x)\,\dd x$, integrating against the density instead of summing against
the pmf. The structure is identical: \emph{apply $g$, weight by how much
probability sits at each input, accumulate}. Sums become integrals, the pmf
becomes the density, but the idea -- average over inputs, not outputs -- is the
same. See the companion calculus reader for why $\int$ is the continuous limit of
$\sum$.
\end{connectionbox}

\begin{example}[LOTUS in action]
\label{ch08:ex-lotus}
Let $X$ be the result of a fair die, and let your payoff be $g(X) = X^2$ dollars.
By LOTUS,
\[
  \E[X^2] = \sum_{x=1}^6 x^2 \cdot \frac16
          = \frac{1 + 4 + 9 + 16 + 25 + 36}{6} = \frac{91}{6} \approx 15.17.
\]
Compare with $(\E X)^2 = 3.5^2 = 12.25$. They differ -- and the gap $15.17 -
12.25 = 2.92$ is precisely the variance of the die, which is the subject of the
next chapter. The strict inequality $\E[X^2] > (\E X)^2$ is a direct case of Jensen's
inequality for the convex function $x^2$, and it is why squaring inflates the
average.
\end{example}

\section{A nonnegative-integer shortcut (the tail-sum formula)}
\label{ch08:tail-sum}

Indicators also give a clean shortcut when $X$ takes nonnegative integer values
$\{0, 1, 2, \dots\}$, expressing its mean directly in terms of \emph{tail
probabilities} -- the chances $\P(X \ge k)$ of reaching or exceeding a threshold $k$:
\[
  \E X = \sum_{k=1}^{\infty} \P(X \ge k).
\]

\begin{whybox}
Why? Use the indicator method on the integer line itself. The clean identity for
a nonnegative integer $X$ is
\[
  X = \sum_{k=1}^{\infty} \ind\{X \ge k\},
\]
which just says ``count up to $X$ by ticking off $1$ for each level $k$ you have
reached.'' (If $X = 3$, the indicators for $k=1,2,3$ are all $1$ and the rest are
$0$, summing to $3$.) Take expectations and swap the (nonnegative) sum with the
expectation:
\[
  \E X = \sum_{k=1}^\infty \E\bigl[\ind\{X\ge k\}\bigr] = \sum_{k=1}^\infty \P(X \ge k).
\]
\end{whybox}

\begin{figure}[H]
\centering
\begin{tikzpicture}[scale=0.95]
  \fill[fill=distblue!20] (0,0) rectangle (8.0, 1.0);
  \fill[fill=sourceblue!25] (2.6,1.0) rectangle (8.0, 2.0);
  \fill[fill=quotientorange!25] (4.8,2.0) rectangle (8.0, 3.0);
  \fill[fill=targetgreen!25] (6.6,3.0) rectangle (8.0, 4.0);

  \draw[thick, draw=black!70] (0,0) -- (8.0,0);
  \draw[thick, draw=black!70] (0,0) -- (0,1.0) -- (2.6,1.0) -- (2.6,2.0) -- (4.8,2.0) -- (4.8,3.0) -- (6.6,3.0) -- (6.6,4.0) -- (8.0,4.0) -- (8.0,0);

  \draw[dashed, draw=gray!60] (2.6,0) -- (2.6,1.0);
  \draw[dashed, draw=gray!60] (4.8,0) -- (4.8,2.0);
  \draw[dashed, draw=gray!60] (6.6,0) -- (6.6,3.0);

  \draw[dashed, draw=gray!60] (0,1.0) -- (2.6,1.0);
  \draw[dashed, draw=gray!60] (2.6,2.0) -- (4.8,2.0);
  \draw[dashed, draw=gray!60] (4.8,3.0) -- (6.6,3.0);

  \node[anchor=north] at (1.3, -0.1) {\footnotesize $\P(X=1)$};
  \node[anchor=north] at (3.7, -0.1) {\footnotesize $\P(X=2)$};
  \node[anchor=north] at (5.7, -0.1) {\footnotesize $\P(X=3)$};
  \node[anchor=north] at (7.3, -0.1) {\footnotesize $\P(X=4)$};

  \draw[decorate,decoration={brace,amplitude=4pt,mirror}] (8.2, 0.05) -- (8.2, 0.95)
    node[midway, right=6pt, font=\footnotesize] {layer $1$: width $=\P(X \ge 1) = 1$};
  \draw[decorate,decoration={brace,amplitude=4pt,mirror}] (8.2, 1.05) -- (8.2, 1.95)
    node[midway, right=6pt, font=\footnotesize] {layer $2$: width $=\P(X \ge 2)$};
  \draw[decorate,decoration={brace,amplitude=4pt,mirror}] (8.2, 2.05) -- (8.2, 2.95)
    node[midway, right=6pt, font=\footnotesize] {layer $3$: width $=\P(X \ge 3)$};
  \draw[decorate,decoration={brace,amplitude=4pt,mirror}] (8.2, 3.05) -- (8.2, 3.95)
    node[midway, right=6pt, font=\footnotesize] {layer $4$: width $=\P(X \ge 4)$};

  \node[align=center, font=\small] at (4.0, 4.6)
    {\textbf{Vertical columns:} $\sum k\,\P(X=k) = \text{Total Area} = \textbf{Horizontal layers:} \sum \P(X \ge k)$};
\end{tikzpicture}
\caption{The tail-sum formula as area: slicing the probability staircase vertically vs.\ horizontally. Summing vertical columns of height $k$ and width $\P(X=k)$ gives the standard definition $\E X = \sum k\,\P(X=k)$. Summing horizontal unit-height slices of width $\P(X \ge k)$ gives the tail sum $\sum_{k=1}^\infty \P(X \ge k)$. Both calculate the exact same total area.}
\label{ch08:fig-tail-sum}
\end{figure}

\begin{example}[Mean of a geometric variable, painlessly]
\label{ch08:ex-geom}
Let $X \sim \Geom(p)$ count the number of trials up to and including the first
success, so $\P(X \ge k) = (1-p)^{k-1}$ (the first $k-1$ trials all failed). Then
\[
  \E X = \sum_{k=1}^\infty (1-p)^{k-1} = \frac{1}{1-(1-p)} = \frac1p,
\]
a clean geometric series, no differentiation needed. If a success has probability
$p = 1/6$ (rolling a six), you expect $6$ rolls to get your first six. Reassuringly
intuitive.
\end{example}

\section{The St.\ Petersburg paradox}
\label{ch08:st-petersburg}

We close with a classic foundational puzzle about expectation -- a problem that
led mathematicians to examine what an expected value actually represents.

\begin{paradoxbox}
\textbf{The game.} A fair coin is flipped until it first comes up heads. If the
first head appears on flip $n$, the casino pays you $2^n$ dollars. So heads on the
first flip pays \$2, heads first on the second flip pays \$4, then \$8, \$16, and
so on, doubling each time you survive another tail.

\textbf{The question.} What is the fair price to play -- how much should you be
willing to pay for one round?

\textbf{Expectation's answer.} The first head lands on flip $n$ with probability
$(1/2)^n$, and that outcome pays $2^n$. So
\[
  \E[\text{payout}] = \sum_{n=1}^{\infty} \underbrace{\left(\tfrac12\right)^{n}}_{\P(\text{first head on } n)}
  \cdot \underbrace{2^{n}}_{\text{payout}}
  = \sum_{n=1}^{\infty} 1 = 1 + 1 + 1 + \cdots = +\infty.
\]
Every term contributes exactly $1$ dollar of expected value, and there are
infinitely many of them. The expected payout is \emph{infinite}. So by the
logic of expected value you should be willing to pay any finite amount -- a
million dollars, a billion -- for a single play. \textbf{Yet essentially nobody
would pay even \$20.} That is the paradox: a mathematically infinite expectation
meeting a flatly finite human willingness to pay.
\end{paradoxbox}

What is going on? Several things at once, and untangling them is genuinely
illuminating.

\begin{happeningbox}
\textbf{First resolution: the infinity is real but lives in tiny tails.} The
expectation is infinite only because of vanishingly rare, astronomically large
payoffs. To even \emph{collect} \$1024 $= 2^{10}$ you need the first head to land
on flip $10$ -- nine tails then a head, probability $2^{-10} = 1/1024$. The median
payout is just \$2. With probability $3/4$ you win \$4 or
less. The mean is dragged to infinity by a sliver of probability promising
unimaginable jackpots. Expectation reports the long-run average exactly; it
just so happens that the long-run average is dominated by events you will almost
never see, and ``almost never'' here means you would have to play for longer than
the age of the universe to make the average behave. The \emph{law of large
numbers fails to kick in on any human timescale}.
\end{happeningbox}

\begin{figure}[H]
\centering
\begin{tikzpicture}
\begin{axis}[
  width=12cm, height=5.6cm,
  title style={font=\small},
  title={St.\ Petersburg Contributions: Constant \$1 Per Term Drives Divergence},
  ybar, bar width=10pt,
  xlabel={flip $n$ of first head}, ylabel={contribution to $\E$ (dollars)},
  xtick={1,2,3,4,5,6,7,8,9,10},
  ymin=0, ymax=1.4, ytick={0,0.5,1},
  enlarge x limits=0.06, axis lines=left,
]
\addplot[fill=trapred!55,draw=trapred] coordinates {
  (1,1)(2,1)(3,1)(4,1)(5,1)(6,1)(7,1)(8,1)(9,1)(10,1)};
\node[anchor=south west,align=left] at (axis cs:1.2,1.02)
  {\small each bar $=\P(n)\cdot 2^n = 1$; the bars never shrink, so the sum diverges};
\end{axis}
\end{tikzpicture}
\caption{St.\ Petersburg: every outcome contributes the same \$1 of expected
value. For typical payout distributions, the tail terms decay rapidly enough for
the sum to converge; here each term is $1$, so the series diverges to $+\infty$.}
\label{ch08:fig-stpete}
\end{figure}

\begin{happeningbox}
\textbf{Second resolution: finite banks cap the game.} An infinite expected value
requires the casino to be able to pay unbounded sums. If the bank holds only
$2^{40}$ dollars (about a trillion) and caps the prize there, the expectation
collapses. Each of the first $40$ outcomes still contributes \$1, and every later
outcome now pays only the capped $2^{40}$, contributing $2^{40}\sum_{n>40}2^{-n}
= 2^{40}\cdot 2^{-40} = 1$ in total:
\[
  \E[\text{capped payout}] = \sum_{n=1}^{40} 1 \;+\; 1 = 41 \text{ dollars}.
\]
A bank of size $2^N$ yields an expected payout of about $N+1$ dollars. Suddenly \$20
is a perfectly sensible price against a realistically finite casino. The infinity
was an artifact of pretending an infinitely rich opponent exists.
\end{happeningbox}

\begin{connectionbox}
\textbf{From coin flips to economics: Daniel Bernoulli and expected utility.}
The resolution published by \textbf{Daniel Bernoulli} in 1738 in the
\emph{Commentaries} of the St.\ Petersburg Academy launched modern decision
theory. He argued that people do not maximize expected \emph{money}; they
maximize expected \emph{utility} -- the subjective personal value or satisfaction
derived from wealth. In particular, money exhibits diminishing marginal utility:
the subjective value of an additional dollar decreases as overall wealth grows.

Bernoulli proposed a logarithmic utility curve, $u(x) = \ln x$. If a player
evaluates the St.\ Petersburg gamble by the utility of the payout $2^n$, the
expected utility is
\[
  \sum_{n=1}^\infty \left(\tfrac12\right)^n \ln(2^n)
  = \ln 2 \sum_{n=1}^\infty \frac{n}{2^n} = 2\ln 2 = \ln 4 < \infty,
\]
using the identity $\sum_{n\ge1} n/2^n = 2$. An expected utility of $2\ln 2 = \ln 4$
equals the utility of a guaranteed \$4 prize. A person with logarithmic utility
would value the gamble at exactly \$4 -- a modest, finite sum that aligns with
human intuition. This 1738 paper laid the foundations for modern economics,
portfolio theory, and the mathematical study of \emph{risk aversion} (the
natural human preference for certainty over a volatile gamble with the same
average payoff).
\end{connectionbox}

\begin{aestheticbox}
The St.\ Petersburg paradox is a lesson about
what expectation \emph{means}. Expectation is the right summary when (a) you can
repeat the gamble many times, (b) the average actually stabilizes on a relevant
timescale, and (c) you care about money linearly. St.\ Petersburg quietly breaks
all three. The takeaway is mature, not cynical: a single number can be exactly
correct and still fail to tell you what to do. Expectation is the beginning of
decision-making under uncertainty, not the end of it -- which is why the next
several chapters add \emph{variance} (how spread out), \emph{concentration} (how
fast averages stabilize), and the laws of large numbers (when, precisely, the
long-run picture is allowed to take over).
\end{aestheticbox}

\section{A gallery of worked examples}
\label{ch08:gallery}

\begin{example}[Expected sum of two dice -- the slow way and the fast way]
\label{ch08:ex-two-dice}
Roll two fair dice; let $S$ be their sum. \emph{Slow way:} write out the
distribution of $S$ (values $2$ through $12$ with the familiar triangular
probabilities $1/36, 2/36, \dots, 6/36, \dots, 1/36$, drawn in
Figure~\ref{ch08:fig-two-dice}) and compute $\sum_s s\,\P(S=s) = 7$. \emph{Fast way:} $S = D_1 + D_2$, and by linearity
$\E S = \E D_1 + \E D_2 = 3.5 + 3.5 = 7$. The fast way needs no distribution at
all, and it generalizes instantly: the expected sum of ten dice is $35$.
\end{example}

\begin{figure}[H]
\centering
\begin{tikzpicture}
\begin{axis}[
  width=12cm, height=5.2cm,
  title style={font=\small},
  title={Dice Sum PMF: Balance Point (Mean) and Peak (Mode) at $S = 7$},
  ybar, bar width=12pt,
  xlabel={sum $s$}, ylabel={$\P(S=s)$},
  xtick={2,3,4,5,6,7,8,9,10,11,12},
  ymin=0, ymax=0.23, ytick={0.05,0.1,0.15},
  yticklabels={$0.05$,$0.10$,$0.15$},
  enlarge x limits=0.06, axis lines=left,
]
\addplot[fill=distblue!55,draw=distblue] coordinates {
  (2,0.0278)(3,0.0556)(4,0.0833)(5,0.1111)(6,0.1389)(7,0.1667)
  (8,0.1389)(9,0.1111)(10,0.0833)(11,0.0556)(12,0.0278)};
\draw[densityred,very thick,->] (axis cs:7,0) -- (axis cs:7,0.185);
\node[densityred,anchor=south] at (axis cs:7,0.185) {\small balance point $=\E S = 7$};
\end{axis}
\end{tikzpicture}
\caption{The triangular pmf of the sum of two fair dice. The peak (mode) and the
balance point (mean) coincide here at $7$ only because the picture is symmetric;
linearity gives $\E S = 7$ with no need to draw the triangle at all.}
\label{ch08:fig-two-dice}
\end{figure}

\begin{example}[Coupon collector: expected time to complete the set]
\label{ch08:ex-coupon}
A cereal company hides one of $n$ equally likely coupons in each box. How many
boxes do you expect to buy to collect all $n$ distinct coupons? Split the wait
into stages: $T_i$ is the number of boxes bought while you hold exactly $i-1$
distinct coupons and are waiting for a new one. When you already have $i-1$
coupons, each box is new with probability $(n-(i-1))/n$, so $T_i$ is geometric
with that success probability and (by Example~\ref{ch08:ex-geom})
\[
  \E T_i = \frac{n}{n-(i-1)}.
\]
The total time is $T = T_1 + \cdots + T_n$, and by linearity
\[
  \E T = \sum_{i=1}^n \frac{n}{n-i+1} = n\sum_{j=1}^n \frac1j = n\,H_n \approx n\ln n.
\]
For $n=50$ coupons, $\E T = 50\,H_{50} \approx 225$ boxes. The stages $T_i$ are
chained together in time, but linearity sums their means without a second thought.
\end{example}

\begin{example}[Expected value of a max via the tail-sum formula]
\label{ch08:ex-max}
Roll a fair die twice and let $M$ be the larger of the two values. Using
$\E M = \sum_{k=1}^6 \P(M \ge k)$ and $\P(M \ge k) = 1 - \P(\text{both} < k) =
1 - \bigl(\tfrac{k-1}{6}\bigr)^2$:
\[
  \E M = \sum_{k=1}^{6}\left[1 - \left(\frac{k-1}{6}\right)^2\right]
       = 6 - \frac{0+1+4+9+16+25}{36} = 6 - \frac{55}{36} = \frac{161}{36} \approx 4.47.
\]
The maximum of two dice averages about $4.47$, comfortably above the single-die
mean of $3.5$ -- taking the better of two tries pulls you upward, as it should.
\end{example}

\begin{walkbox}
You are about to flip a fair coin $100$ times. A friend bets you that the number
of times the pattern \texttt{HH} appears (in overlapping windows -- positions
$1\text{-}2, 2\text{-}3, \dots, 99\text{-}100$) will be more than $30$. Should you
take the bet? Use indicators: there are $99$ overlapping windows, and each window
shows \texttt{HH} with probability $\tfrac14$. By linearity the expected count is
$99\cdot\tfrac14 = 24.75$, regardless of the heavy overlap between consecutive
windows. Since $24.75 < 30$, betting against ``more than $30$'' is the favorable
side -- and again, the dependence among overlapping windows never entered the
mean calculation.
\end{walkbox}

\section{What to carry forward}
\label{ch08:summary}

\begin{focusbox}
\textbf{The five things that matter.}
\begin{enumerate}
  \item \textbf{Two pictures.} $\E X = \sum_x x\,\P(X=x)$ is both a long-run
  average and a center of mass. It need not be a possible value.
  \item \textbf{Indicators bridge to probability.} $\E[\ind_A] = \P(A)$. Every
  probability is the expected value of an indicator.
  \item \textbf{Linearity is unconditional.} $\E[\sum_i X_i] = \sum_i \E X_i$
  with \emph{no} independence assumption -- the most useful tool in the book.
  \item \textbf{The indicator method.} Write a count as $\sum_i \ind_{A_i}$,
  then $\E = \sum_i \P(A_i)$. It cracked hats, records, coupons, QuickSort.
  \item \textbf{LOTUS.} $\E[g(X)] = \sum_x g(x)\,\P(X=x)$: average over inputs,
  weighting by the original pmf, never deriving the distribution of $g(X)$.
\end{enumerate}
One caution: an infinite or heavy-tailed expectation (as in the St.\ Petersburg
game) can be mathematically exact yet practically misleading. The mean is a single summary number,
and a summary inevitably leaves things behind. Two games can share an identical
expectation of \$50 -- one paying a crisp guaranteed \$50 bill, the other paying
\$0 or \$100 on a coin flip -- yet represent completely different realities of
risk. Expectation locates the fulcrum where the distribution balances; it says
nothing about how widely the probability mass is scattered away from that center.
To measure spread, risk, and stability, we need to average the squared deviations
from the mean -- the \emph{variance} -- and examine the higher \emph{moments} that
capture asymmetry and heavy tails. That is the subject of Chapter~9.
\end{focusbox}

\section*{Exercises}
\addcontentsline{toc}{section}{Exercises}

\begin{exercise}[Warm-up means]
\label{ch08:exr-warmup}
Compute $\E X$ directly from the definition for each:
(a) $X$ is the number on a fair four-sided die ($1,2,3,4$);
(b) $X$ pays \$10 with probability $0.1$ and \$0 otherwise;
(c) $X \sim \Bern(p)$.
\end{exercise}

\begin{exercise}[Center of mass is not the mode]
\label{ch08:exr-com}
A random variable takes value $0$ with probability $0.9$ and value $100$ with
probability $0.1$. Find $\E X$. Explain in one sentence why the expectation lies
in a region the variable visits only $10\%$ of the time, using the center-of-mass
picture.
\end{exercise}

\begin{exercise}[Linearity, no independence]
\label{ch08:exr-linearity}
Draw $3$ cards without replacement from a standard $52$-card deck. Let $A$ be the
number of aces drawn. Using indicators $A = A_1 + A_2 + A_3$ (one per drawn
position), find $\E A$. Then explain why the same answer holds even though the
three draws are dependent.
\end{exercise}

\begin{exercise}[Hat-matching variations]
\label{ch08:exr-hats}
$n$ couples (so $2n$ people) are seated at a round table by a random permutation.
(a) What is the expected number of couples seated next to each other?
\emph{Hint:} indicator per person for ``my right neighbor is my partner.''
(b) Why is the dependence between these indicators irrelevant to the mean?
\end{exercise}

\begin{exercise}[Fixed points of a random function]
\label{ch08:exr-fixedpoints}
A function $f:\{1,\dots,n\}\to\{1,\dots,n\}$ is chosen so that each $f(i)$ is an
independent uniform value in $\{1,\dots,n\}$ (note: not necessarily a
permutation). Let $X = \#\{i : f(i) = i\}$ be the number of fixed points. Find
$\E X$. Compare with the permutation (hat) case and comment.
\end{exercise}

\begin{exercise}[LOTUS practice]
\label{ch08:exr-lotus}
Let $X$ be uniform on $\{1,2,3,4,5\}$.
(a) Compute $\E\!\left[\tfrac{1}{X}\right]$ using LOTUS.
(b) Compute $\E[(X-3)^2]$.
(c) Verify that $\E[(X-3)^2] \ne (\E X - 3)^2$ and explain which inequality
(Jensen) is at work.
\end{exercise}

\begin{exercise}[Tail-sum formula]
\label{ch08:exr-tailsum}
(a) Use $\E X = \sum_{k\ge1}\P(X\ge k)$ to find the expected number of rolls of a
fair die needed to first see a number $\ge 5$.
(b) Let $X$ be the minimum of three independent fair die rolls. Using
$\P(X \ge k) = \bigl(\tfrac{7-k}{6}\bigr)^3$ (all three rolls must be $\ge k$) for
$k = 1,\dots,6$, compute $\E X$. \emph{Answer:} $\E X = \sum_{k=1}^{6}\bigl(\tfrac{7-k}{6}\bigr)^3
= \tfrac{6^3+5^3+4^3+3^3+2^3+1^3}{216} = \tfrac{441}{216} \approx 2.04$.
\end{exercise}

\begin{exercise}[St.\ Petersburg with a cap]
\label{ch08:exr-stpete}
Modify the St.\ Petersburg game so the casino pays $\min(2^n, 2^{20})$ -- the prize
is capped at about a million dollars.
(a) Show the expected payout is finite and find it (approximately).
(b) Explain in words how a finite bank turns an infinite expectation into a small
one.
\end{exercise}

\begin{exercise}[Proof: linearity from first principles]
\label{ch08:exr-proof-linearity}
Let $X$ and $Y$ be discrete random variables on a finite sample space $\Omega$
with outcome probabilities $\P(\omega)$. Starting from $\E Z = \sum_{\omega \in
\Omega} Z(\omega)\,\P(\omega)$, prove that $\E[aX + bY] = a\,\E X + b\,\E Y$ for
constants $a, b$. Identify the single step where you use the fact that we are
adding numbers \emph{pointwise at each outcome}, and explain why no assumption
about independence is needed.
\end{exercise}

\begin{exercise}[Proof: LOTUS]
\label{ch08:exr-proof-lotus}
Prove the discrete LOTUS, $\E[g(X)] = \sum_x g(x)\,\P(X=x)$, by grouping the
$x$-values according to the value $y = g(x)$ they map to. Where in the argument
does the possibility that $g$ is \emph{not} one-to-one get correctly handled?
\end{exercise}

\begin{exercise}[Indicator method in the wild]
\label{ch08:exr-balls}
You throw $n$ balls independently and uniformly into $n$ bins.
(a) Find the expected number of empty bins (exactly, then as an approximation for
large $n$).
(b) Find the expected number of bins containing exactly one ball.
(c) These two indicator families overlap and are dependent. State precisely why
that does not affect either expected count.
\end{exercise}

\begin{exercise}[Conceptual: when is the mean a bad summary?]
\label{ch08:exr-conceptual}
Give two qualitatively different real-world situations in which the expectation of
a quantity is a misleading thing to report, and say what you would report instead.
One of your examples should involve a heavy/long tail (St.\ Petersburg flavored)
and the other should involve a strongly bimodal distribution (one with two distinct
peaks, like camel humps, rather than a single cluster around the center).
\end{exercise}
